\documentclass[conference]{IEEEtran}
\usepackage{booktabs}
\usepackage{url}
\usepackage{listings}

% Milestone mini-report: Benchmark v4 (executable-checker benchmark)
% Two-page target. Numbers below are generated by the build itself:
% exports/gate_report.json, exports/tier_report.json, benchmark_v4_decisions.md.
% Style: plain prose, no em-dashes, honest representation of failures.

\begin{document}

\title{Milestone Report: An Executable-Checker Benchmark for CHI@Edge Coding Assistance (Benchmark v4)}

\author{\IEEEauthorblockN{Vivek Rai}
\IEEEauthorblockA{OSRE 2026 Fellowship, UC Santa Cruz OSPO Summer of Reproducibility \\
Project: From Lookup to Reasoning, Chameleon Cloud CHI@Edge \\
Mentor: Marc Richardson}}

\maketitle

\begin{abstract}
Milestone 1 established that grounding language models in Trovi reproducibility
artifacts lifts CHI@Edge coding accuracy from under 22\% to above 82\%
regardless of model, but that benchmark relied on a single human grader and 28
prompts. This milestone replaces the measurement instrument. Benchmark v4
contains 50 items whose ground truth is machine-verifiable: every item ships
AST-based checkers, and a gold answer is admitted only after passing 100\% of
its own checkers, following the SWE-bench principle that ground truth must be
executable. The gate admits 50 of 50 golds; a ten-control negative suite
confirms the checkers reject bare-metal substitutions, fabricated device
profiles, the filter\_reserved direction bug, invented lease-extension calls,
and Kubernetes-unsafe names. Difficulty tiers are computed from key-token
coverage of the fed context rather than asserted, and every check carries a
mechanism, specifics, or safety tag, which turns the benchmark into the
measurement instrument for the project's generalization hypothesis.
\end{abstract}

\section{Context and motivation}
The project extends the Chameleon documentation chatbot into an
allocation-aware resource advisor that reads live CHI@Edge state, reasons about
a plain-language workload, validates the result, and emits a runnable
python-chi reservation. Milestone 1 measured the baseline with a 28-prompt
benchmark graded by hand on a 0 to 2 rubric \cite{m1}. Two weaknesses were
identified in planning review: a single unblinded grader is a validity risk,
and prose golds cannot be re-verified mechanically as the python-chi API
evolves. Benchmark v4 addresses both before any system-level evaluation of the
advisor is run, so that the yardstick is fixed first.

\section{Source material}
Six public artifacts were extracted with a structured, grep-verified protocol
(commit-pinned; every API call, keyword argument, and magic string confirmed
verbatim against source): the SSH tutorial (A1), Pi Camera (A2), Sense HAT in
two variants (A3), the OpenMCP edge LLM notebooks (A4), serve-edge-chi on
Raspberry Pi 5 (A5), and edge-cpu-inference (A6). A1 and A6 use the imperative
python-chi generation; A2 through A5 use the object-oriented generation. A4
deliberately retains its KVM@TACC bare-metal notebooks in the grounding file as
distractor content, so feeding A4 also tests whether a model drifts to
server-style APIs when asked about edge (decision D08).

\section{Benchmark design}
\textbf{Items.} 50 items: the 28 v3 prompts ported verbatim for continuity and
calibration (D11), 18 new items covering parameter adaptation, cross-artifact
composition, GPU runtime, API-generation translation in both directions, and
deliberately uncovered requests (microphone, UVC webcam, lease extension), and
4 availability items that execute against a recorded inventory snapshot.

\textbf{Checkers.} Each item carries 1 to 14 checks from a shared AST library:
required and forbidden calls, keyword values with cross-generation aliases
(machine\_type or machine\_name, amount or count, image\_ref or image),
reservation wiring, lease duration, container-name constraints, availability
query direction, a device-profile allowlist, and abstention-or-discovery
detection. Forbidden-call checks operate on the AST only and string checks on
code blocks only, so prose that warns against an API is never penalized (D10).
Golds are authored in the object-oriented generation; imperative answers earn
full credit iff trap-free, including platform\_version=2 on create\_container
(D01).

\textbf{Admission gate.} A gold enters the bank only after passing all of its
own checkers, and availability golds must additionally reproduce expected
output when executed against the snapshot through a stub chi package. The gate
admits 50 of 50. Because a gate that never fails is itself untested, a
ten-control negative suite was run: deliberately wrong answers for the known
failure modes must be rejected, a prose mention of a forbidden call must not
be, and the snapshot execution path must demonstrably fire. All ten controls
behaved as designed (D17).

\textbf{Computed tiers.} Difficulty tiers are a property of an item and a fed
context, computed from key-token coverage: T1 direct, T2 adapted values, T3
multi-artifact composition, T4 uncovered. Across 114 item-condition pairs the
matched condition yields 32 T1, 7 T2, 6 T3, and 2 T4; the 14-item held-out
subset is uniformly T4 by construction. The two matched-condition T4 items are
themselves a finding: no artifact documents the supported\_device\_profiles
field, so the correct answer to profile-discovery questions is outside every
tutorial (D15).

\textbf{Hypothesis instrument.} Every check is tagged mechanism, specifics, or
safety, and results are aggregated per tag per condition (D03). Under held-out
feeding, the transfer hypothesis H1 predicts mechanism checks pass while
specifics checks fail; the lookup hypothesis H0 predicts both fail. The
benchmark therefore measures the generalization question directly instead of
inferring it from aggregate scores.

\section{Failures and reroutes}
Three are worth recording. First, the tier assigner flagged its own author:
P28 carried the generation-specific token container.Container, absent from the
imperative A1 grounding, and was computed T4 against an intended T1. The token
policy was corrected and the pair re-verified (D16). Second, one negative
control was initially over-specified, expecting a redundant second check to
fire on the platform\_version omission; the answer was correctly rejected and
the test assertion, not the item, was fixed (D17). Third, duration checks fail
when a duration is not statically determinable, for example manual start and
end date strings; this is a known false-negative surface that requires hand
adjudication and is documented rather than papered over (D12).

\section{Limitations}
Live testbed execution of golds (V2) was dropped for this cycle by decision of
the fellow; no lease is ever created by the benchmark (D05). The shipped
snapshot is synthetic and clearly flagged; a recorded Blazar snapshot from the
replay harness must replace it before availability numbers are reported (D13).
Scoring executes model code in a subprocess for availability items and should
run in a disposable environment. Checker-versus-human agreement is not yet
measured; the 168 graded v3 cells share prompts with the ported subset and
will serve as the calibration set.

\section{Next steps}
(1) Run the checkers over the 168 graded v3 cells and report agreement with
the human grades. (2) Replace the synthetic snapshot with a recorded one.
(3) Run the frozen benchmark subjects and the chi-edge-advisor itself over all
50 items in the blind, matched, and held-out conditions, reporting advisor
results with its validator on and off (D07). (4) Fold the per-group results
into the generalization section of the workshop paper.

\begin{thebibliography}{9}
\bibitem{m1} V. Rai, ``Milestone 1: A blind versus artifact-grounded coding
benchmark for CHI@Edge,'' OSRE 2026 project report, 2026.
\bibitem{swebench} C. Jimenez et al., ``SWE-bench: Can Language Models Resolve
Real-World GitHub Issues?'' ICLR, 2024.
\bibitem{chameleon} K. Keahey et al., ``Lessons Learned from the Chameleon
Testbed,'' USENIX ATC, 2020.
\bibitem{pychi} Chameleon Cloud, ``python-chi,''
\url{https://github.com/ChameleonCloud/python-chi}.
\bibitem{artifacts} Artifact repositories A1--A6, commit-pinned in
\texttt{extractions/}: edge\_ssh\_image, edge-picamera-image,
edge\_sensehat\_image, MCP\_SLM\_Project, serve-edge-chi, edge-cpu-inference.
\end{thebibliography}

\end{document}
